% Agrippa's trilemma, counted — preprint for Zenodo. Build: lualatex twice
% (FreeSerif/FreeMono are needed for ¬ ∧ ∨ ⊕ ↔). Russian source: AGRIPPA-RU.md.
\documentclass[11pt,a4paper]{article}
\usepackage{fontspec}
\setmainfont{FreeSerif}[RawFeature={-ccmp;-mark;-mkmk}]
\setmonofont{FreeMono}
\usepackage{amsmath,amssymb}
\usepackage{geometry}
\geometry{margin=1in}
\usepackage{enumitem}
\setlist{nosep,leftmargin=1.4em}
\usepackage[hidelinks,pdftitle={Agrippa's Trilemma, Counted: Acts, Posits and the Formal Register},pdfauthor={Vitaly Reznik}]{hyperref}
\usepackage{microtype}
\setlength{\emergencystretch}{3em}

\newcommand{\T}{\mathsf{T}}
\newcommand{\F}{\mathsf{F}}
\newcommand{\Z}{\mathsf{Z}}
\newcommand{\ZTL}{\textsc{ztl}}
\newcommand{\code}[1]{\texttt{#1}}

\title{\textbf{Agrippa's Trilemma, Counted:\\ Acts, Posits and the Formal Register}\\[6pt]
\large A kernel-checked reading of the regress, the circle and the dogma}
\author{Vitaly Reznik\\[2pt]
\small Independent researcher \quad ORCID 0009-0002-4103-6387 \quad \texttt{vitalyreznik@gmail.com}}
\date{\small October 2026 \quad---\quad Lean 4 sources: \url{https://github.com/inventor1975/ZTL}}

\begin{document}
\maketitle

\begin{abstract}
\noindent
Agrippa's trilemma says that a justification either regresses without end, closes into a circle, or breaks off at a
dogma. We read it by one number: how many admissible readings a chain of reasons has, under classical (supported)
semantics. On a finite base the regress cannot occur (on an infinite one it reads twice, like the dogma). An odd loop
has no reading: it is a paradox. An even loop and the dogma ``$p$ because $p$'' have two
each: they are one and the same fork. Only a chain that reaches a ground --- a reason that rests on nothing --- reads once. Grounds come in two kinds. The
value of an \emph{act} is computed from its content and is the same under every base. The value of a \emph{posit}
(an axiom, a Thing) is set by the base, and a rival base always exists. We read this as the thesis that knowledge rests on acts alone, while about Things Agrippa is right and his horns are
the signature of the formal register. The formal results are proved in Lean~4 on the empty axiom list (98 theorems in four modules, auxiliary lemmas included); the thesis is our reading of them, and where it goes beyond them we say so. The count itself is old (Kripke's grounded, ungrounded and paradoxical sentences); we report the literature
found and read, and claim novelty only for applying it to Agrippa, for the act/posit line and for the machine-checked
proofs.
\end{abstract}

\section{The problem}

Sextus transmits Agrippa's five modes: disagreement, regress, relativity, hypothesis, circularity (\emph{Outlines of
Pyrrhonism} I.164--177). Albert reduced three of them to the ``Münchhausen trilemma'' (1968): every justification ends
in a regress, a circle or a dogmatic break-off, so nothing is justified.

We do not dispute the trilemma's logic; its case analysis is sound. We dispute that its three outcomes are of equal
standing. The question we ask of each is: how many values are compatible with the chain of reasons?

\section{The model}

A justification system on a base of $n$ statements. Each statement rests either on a \emph{nullary ground} (its value
is given and nothing further is asked) or on another statement $j$ --- on $j$ itself or on its negation. The base is
\emph{closed}: the reasons of the first $n$ statements are among them.

A \emph{reading} of a statement $x$ is an assignment of truth values under which every support on $x$'s chain holds.
Only the supports of that chain count; the rest of the system does not affect the count.

The semantics is classical (supported): ``$p$ because $p$'' admits both true and false. This is a choice, and the
count depends on it. In logic programming the rule $p \leftarrow p$ has one stable model (the empty one): minimality
chooses false, and under that semantics the dogma would read once. Clark's completion $p \leftrightarrow p$ has two
models; for tight programs (no positive loops) the two semantics coincide (Fages 1994). We do not adopt minimality
because it settles the dogma by a default --- false unless supported --- that the chain itself does not supply; that
default is the kind of stop we are examining. A reader who accepts minimality as a neutral principle will count
differently, and not only the dogma: the positive circle $a := b,\ b := a$ then also reads once (the empty model),
while the circle with two negations still reads twice, so under minimality ``the dogma and the even circle take the
same place'' no longer holds.

Finiteness matters for the constructive proofs, not for the count: an infinite chain of distinct reasons reads
twice (\code{regress\_two} below), the same fork as the dogma. Yablo's paradox (1993) needs a statement resting on
infinitely many others at once, which is outside both models here (a reason is one statement, or a finite
expression).

\section{The count}

Module \code{lean/Agrippa.lean}.
\begin{itemize}
\item \code{no\_regress}: on a closed base of $n$ statements a chain repeats within $n$ steps (pigeonhole). An
  infinite regress of distinct reasons needs an infinite base.
\item \code{regress\_two}: on an infinite base the chain ``$n$ because $n + 1$'' has two readings (all true, all
  false). The regress horn, where it can occur, is the fork.
\item \code{chain\_value}: along a reading, the value of $x$ is the parity of negations on the path combined with the
  value reached.
\item \code{odd\_zero}: a loop with an odd return has no reading. Example: the liar $x := \neg x$ (\code{liar\_zero}).
\item \code{even\_two}: a loop all of whose returns are even has two readings that disagree on $x$.
\item \code{forced\_one}: a chain that reaches a ground has one reading.
\item \code{dogma\_two}, \code{circle\_two}, \code{twoneg\_two}: the dogma $f := f$, the circle $a := b,\ b := a$ and
  the circle $a := \neg b,\ b := \neg a$ read twice. \emph{The dogma and the even circle take the same place in the
  count.}
\item \code{knowledge\_iff\_ground}: $x$ has exactly one reading if and only if its chain reaches a ground.
\end{itemize}
The last theorem is proved without excluded middle: the parity of returns is decided by a bounded search over one
turn of the loop (\code{oddBelow}) and carried to every return (\code{all\_even}).

In sum: on a finite base the regress is impossible; where it occurs it gives 2; the circle gives 0 or 2; the dogma
gives 2; a chain reaching a ground gives 1. For chains, one count sorts every horn. Classically
\code{knowledge\_iff\_ground} also holds on an infinite base; finiteness is what makes the proof constructive. Klein's
infinitism is not refuted: an infinite regress is admitted, and counted as a fork.

\section{Several reasons at once}

A chain is a simplification: a statement rests on one reason. Module \code{lean/AgrippaNet.lean} lets a statement
rest on an expression over other statements built with ``and'', ``or'', ``not''.
\begin{itemize}
\item \code{knows\_of\_good} (sufficiency): if the reasons bottom out, in finitely many steps, in acts alone, a reading
  exists under every base and every reading under every base gives one value.
\item \code{knows\_unless\_loop\_or\_posit} (necessity, constructively): on a finite closed base, the absence of loops
  and posits among the reasons implies knowledge. By contraposition (classically; constructively one gets
  its double negation), whatever is not knowledge rests on a loop (circle, dogma, liar) or on a posit.
\item \code{act\_or\_dogma} (counterexample): $x := (\text{the act } 0 = 0) \lor (\text{a dogma})$ has one value under
  every base, although a dogma is among its reasons. A true act overrides the dogma here; with a false act ($0 = 1$)
  the dogma decides and $x$ reads twice (\code{false\_act\_or\_dogma\_two}). The right formulation is therefore ``the value is forced by acts'',
  not ``no dogma among the reasons'' --- and with several reasons the count classifies statements, not horns.
\item \code{dogma\_net\_two}, \code{liar\_net\_zero}: the controls in the richer language.
\end{itemize}
The connectives are given as the medium; they are not derived from the act.

\section{Act and posit}

The count separates a ground from a self-reference, not an act from an axiom: an axiom written as a ground also reads
once. We draw the line with the two registers of VR-Forms (Reznik 2026b): an \emph{act} is constructed, a \emph{Thing}
is posited, and every axiom is optional --- a system may reject it and go on reasoning. (The act/Thing names are
introduced in VR-Forms version 1.0.3, §I.3.)

In Lean a ground is of two kinds. An act is a claim $m = n$ whose value is computed (\code{Nat.beq}). A posit is a value
set by the base $\beta$.
\begin{itemize}
\item \code{act\_absolute}, \code{act\_reads}: if $x$'s chain reaches an act, a reading exists under every base,
  and all readings under all bases give $x$ one value.
\item \code{posit\_relative}, \code{posit\_each\_once}: if the chain reaches a posit, each base gives exactly one
  reading, and two bases give $x$ opposite values.
\end{itemize}

The act is a bounded computation over given numerals. By Hamkins (2012), different models of set theory have
different natural numbers and different arithmetic truths. This does not reach our acts: an equality $m = n$ of
numerals is proved or refuted already in Robinson's arithmetic Q, so it has the same value in every model of
arithmetic, nonstandard ones included. Statements with unbounded quantifiers over all numbers --- for instance the
consistency of ZFC --- are not acts in our sense and may themselves belong to the formal register. We do not claim
that all of arithmetic is settled.

Both theorems nearly follow from the definitions: the posit is a free parameter, the act a computed value. The
substantive step lies outside them, in classifying a real statement. That the axiom of choice and the continuum
hypothesis are posits follows from independence (Gödel 1938, Cohen 1963), not from these theorems.

\section{In \ZTL's own terms}

\ZTL\ is a three-valued logic whose connectives are two-valued: the value $\Z$ occurs only on atoms, and a verdict is
always $\T$ or $\F$. $\Z$ is a \emph{dirty} parameter: it is not ours, another system controls it, and it may change
between two reads. A connective therefore reads $\Z$ at each occurrence separately --- both readings, independently ---
and returns $\T$ only if $\T$ is forced; so $p \to p$ with $p = \Z$ is not a law (one occurrence may read true, the
other false). The word ``unverified'' would mislead here: it suggests one value that is merely unknown, which is the
supervaluation reading \ZTL\ rejects. Verification is not a property of $\Z$ but a way to clean it: the dirty
parameter is replaced by a checked value, a new and clean input, rather than its hidden value being disclosed. (The stand \code{dilemmas/dirty\_param.py} separates the substituted, the
stable and the foreign reading and shows that \ZTL\ reads the third.)

The number of readings shows not in the verdict but in its future. Module \code{lean/AgrippaZ.lean}: an admissible
future of $x$ is a cleaning of every $\Z$ to $\T$ or $\F$ under which the chain's supports hold. Quantifying over
cleanings is not supervaluation: today's verdict is still computed occurrence by occurrence; the futures only grade
it. Both sides are compared in the same state, before anything is cleaned.
\begin{itemize}
\item \code{forced\_futures\_agree}: a chain reaching a ground has an admissible future, and every admissible future
  gives the same verdict. Knowledge earned.
\item \code{two\_readings\_futures\_split}: with two readings, two admissible futures give opposite verdicts.
  Knowledge on credit.
\item \code{paradox\_no\_future}: no reading, no admissible future.
\end{itemize}

\emph{The descent of a Thing} (module \code{lean/ZDescent.lean}). A formal atom is replaced by its operational meaning
($\pi$).
\begin{itemize}
\item \code{conservative}: if $\pi$ of every formal axiom follows from the operational theory, the formal register adds
  no operational consequence.
\end{itemize}
The consequence relation is semantic (\ZTL\ valuations); no proof calculus is used, and a substitution lemma
suffices (\code{descend\_eval}). The premise --- every formal axiom is operationally harmless --- is the substantive
condition, as in Hilbert; Shapiro (1983) shows through Gödel sentences that it fails for real theories that prove new
consistency statements, so it must be checked Thing by Thing. VR-Forms (Reznik 2026b) proves the same shape
for a classical system with Peirce's law. Classical logic is the case of \ZTL\ in which every atom is settled, so the
classical proof does not carry over by itself.
\begin{itemize}
\item \code{every\_base\_is\_not\_knowledge}: ``$A$ or not $A$'' is true in the base with $A$ and in the base with
  $\neg A$, but does not follow in \ZTL. A base fixes $A$; a dirty parameter is not fixed by anything until it is
cleaned. The operational part of a
  Thing is what is constructed without it, not what is ``true in every base'' (supervaluation).
\end{itemize}

\section{Reading Agrippa}

Sextus calls dogma assent to the non-evident, to how things are by nature (I.13). The sceptic does not abolish
appearances: doubt concerns the account given of an appearance, not the appearance (I.19--20). He lives by nature,
the passions, laws and customs, and instruction in the arts (I.21--24). Against hypothesis Sextus sets: the opposite
hypothesis deserves no less credence (I.173--174) --- close to what we call a rival base.

Two sides follow.
\begin{itemize}
\item \emph{The trilemma does not reach acts.} An act rests on nothing, reads once, and no base changes it. Sextus
  aims the modes at assent, not at acting (I.19--24). His ``acting'' is wider than our technical act: laws, customs
  and the arts are, by our own definition, posits; the sceptic follows them without assent, which is the point.
\item \emph{About Things Agrippa is right.} A posit has a rival base: Euclidean and Lobachevskian geometry, ZF with
  choice and with its negation (both proved consistent relative to ZF); a world with centaurs and without (an
  informal case). The modes of disagreement and relativity
  describe exactly this. Agrippa's horns are the signature of the formal register.
\end{itemize}

\emph{Empirical knowledge.} ``It is raining'' is a posit: the world sets its value, and its twin is thinkable. It is
knowledge on credit. Earned knowledge is only what acts force. Reading an observation as an act of measurement moves
the credit onto the instrument, whose reliability is itself a premise.

\emph{The thesis, stated.} In this paper ``knowledge'' is a technical term: exactly one reading, the same under
every base. The philosophical claim is that knowledge in this sense \emph{rests} on acts. Whether the ground itself
deserves the word is a decision about the word, not a theorem; Wittgenstein places at the bottom an acting that is
not itself known (\emph{On Certainty}, §§204, 253).

\emph{The paradox} is not a third entity beside the act and the Thing. It has no reading, no model, no future. We
read it as a failed claim to become a Thing: a posit that has no model to be admitted by.

\section{The critics}

The stand \code{dilemmas/agrippa\_critics.py} runs 17 objections through \ZTL's judge (\code{ztljudge}). For each: the
logic of the inference (conceded in full if it holds at every value), then the premises under values taken \emph{from
computation} --- the count of readings on small systems --- not written in by hand. A premise nothing computes stays
dirty ($\Z$), and the instrument says whether it can ever be cleaned.

The objections and their formulations are ours, and so is the judge; the stand is a disciplined check, not an
independent verdict. Result (17 rows): 10 objections, as formalized here, fall on a premise whose value the count computes as false in the
model; 2 stand, both controls built to test the same argument against the dogma and the posit; 5 rows rest on credit
(Klein and Peirce share one row) --- four of them permanent, which neither side can repay, and one, Kant's,
redeemable in principle by exhibiting the table of categories as forced, which is not done.
\begin{itemize}
\item \emph{Fall (as formalized):} Albert against the act, Popper, Quine, Albert's second move (``who chose the base?'') against the
  act, Sellars/Davidson within the model, BonJour, Hegel, foundationalism, Descartes, Fichte. Fichte falls
  instructively: the self-positing ``I = I'' is ``$p$ because $p$'', two readings. The primary act rests on nothing,
  not on itself. Descartes falls on the circle of the guarantee (clear perception $\leftarrow$ God $\leftarrow$ clear
  perception; Arnauld's Fourth Objections), not on the cogito.
\item \emph{Stand (controls):} Albert against the dogma and Albert against a posit. The argument that fails against the act
  holds against the dogma and against the posit.
\item \emph{Credit:} Kant (the given table of categories), Sellars (does a real act yield a value), Sextus (is the act
  a mere appearance), Klein and Peirce (an infinite base), Wittgenstein (the word ``knowledge'').
\end{itemize}

\section{Prior work and the limits of our search}

The literature check (about thirty English-language searches, October 2026) was carried out by AI-assisted
sessions (see the disclosure) and then checked by an independent AI review, which caught two errors, corrected
here. Below, ``read'' means read in full by those sessions; the author has not re-read every source. Not finding a work
does not mean it does not exist: an author unknown to us may have done the same.

\emph{Known before us, and read:}
\begin{itemize}
\item The count as a classifier. Kripke (1975) already separates grounded sentences, the ungrounded but consistent
  truth-teller and the paradoxical liar --- our 1, 2 and 0; see also Herzberger (1970), Yablo (1982), Leitgeb (2005) on
  grounding and dependence, and Richardson (1953) on kernels of digraphs. In Gaifman (1992) the truth-teller has two
  self-supporting valuations and the liar none. Walicki and Dyrkolbotn (2022) use the graph normal form, in which our structure of supports appears: a ground is
  a sink, models are kernels, the liar has none.
\item The regress on graphs. Miao (2023) and Berker (2015) describe the trilemma by a root, a cycle and an infinite
  path, without counting values.
\item Splitting the third horn into a break-off and an ungrounded break-off: Day and Athos (2026). They have no count,
  and their reliable stop needs warrant from outside the chain; ours is forced within the model. Their work is
  independent of ours (self-published; one author is an AI persona).
\item Independence as a mark of the posited: Hamkins (2012, the multiverse) and Feferman (2011, the indefiniteness of
  CH). Feferman's criterion is the provability of excluded middle, not a rival model. He holds CH to be indefinite:
  formally definite in his semi-constructive SCS with full power set, but, as he conjectured and Rathjen (2016)
  proved, not definite with the power set of $\omega$ only.
\item The act as ground: Fichte, Hintikka (1962), Wittgenstein. None of them has invariance under the base.
\item Conservativity as the criterion for admitting the ideal: Hilbert (1926), Field (1980). Shapiro's (1983)
  objection through Gödel sentences applies to our descent too.
\item The fork of two equally coherent systems is a known objection to coherentism.
\end{itemize}

\emph{Seen only in abstract or summary:} Walicki, ``Reference, paradoxes and truth'' (2009, closed access); Kripke
(1975); Gupta and Belnap (1993); Barnes (1990), a monograph on the Agrippan modes themselves --- a gap we name; Dung
(1995); Fages (1994); Cook (2004).

\emph{Not found in anything we read:}
\begin{itemize}
\item the 0 / 1 / 2 count applied to Agrippa's horns (the count itself is Kripke's);
\item the constructive equivalence ``exactly one reading iff the chain reaches a ground'' (the kernel literature
  above gives existence and counts, not this equivalence for chains) and its form with and/or/not;
\item the distinction ``an act is base-invariant / a posit has a rival base'' as theorems;
\item ``two readings = on credit'' in a three-valued logic;
\item the contrast between ``constructed without the Thing'' and ``true in every base''.
\end{itemize}
We claim novelty only for this combination and for the machine check. The parts are known mathematics and known
philosophical moves.

\section{Limitations}

\begin{itemize}
\item The base is finite in the constructive results; an infinite regress is counted (two readings), but Yablo-style
  reasons resting on infinitely many statements are outside the models.
\item The model is two-valued in its readings. The link to three-valued \ZTL\ is given in Section~6 through the
  verdict's future: admissible futures are cleanings of dirty atoms, not their hidden values.
\item The act is a decidable equality of numerals --- a thin base for a general epistemology. Its absoluteness is
  relative to the medium: the logic, the connectives and the numerals the chain is read in. An attack on the medium
  itself is not closed by a theorem.
\item The Lean kernel checks that the proofs prove the stated theorems. It does not check that the definitions
  (reading, act, posit, base) model Agrippa, Sextus or Hamkins faithfully; that is for the reader to judge from the
  definitions, which are short.
\item Fogelin (1994) makes ordinary knowledge relative to a level of scrutiny; whether our act survives a stricter
  standard than the model's is not shown.
\item Classifying real statements into acts and posits is done outside the model.
\item Whether the Pyrrhonist has any beliefs at all, everyday ones included, is disputed (Burnyeat 1980, Barnes
  1982 against Frede 1979). Our reading, that the modes target assent and not acting, rests on I.19--24 and sides with
  Frede.
\end{itemize}

\section*{Reproducibility}

All theorems are in \url{https://github.com/inventor1975/ZTL}, directory \code{lean/}, at the tag
\code{agrippa-preprint-v1.0} (commit f8c4fae): \code{Agrippa} (62 theorems), \code{AgrippaNet} (24),
\code{AgrippaZ} (8), \code{ZDescent} (4) --- 98 in all, auxiliary lemmas included. History: 55857aa (the count),
5ab548f (\code{knowledge\_iff\_ground}), 0e6b3f2 (acts and posits, the critics' stand), fbcd9b2 (\code{ZDescent}),
66fcf06 (\code{AgrippaNet}), bc51055 and f8c4fae (\code{AgrippaZ} and the corrections after review). Each theorem is
on the empty axiom list; to check: \code{cd lean \&\& lake build}, then \code{python3 inventory/axiom\_audit.py} from
the repository root, which audits every theorem individually. The critics' stand: \code{python3
dilemmas/agrippa\_critics.py}; the three readings of a dirty parameter: \code{python3 dilemmas/dirty\_param.py}.

\section*{Acknowledgements and AI disclosure}

The text, the Lean proofs and the stand's code were prepared with the substantial participation of the AI system
Claude (Anthropic), model Opus~5.5, in a dialogue setting: the system drafted the text, the proofs and the code under the author's direction. Every decision
on the question, the frame and the wording was made by the author, Vitaly Reznik, who bears the responsibility for
the content. The literature check was carried out by auxiliary sessions of the same system. In accordance with COPE
and ICMJE recommendations the AI system is not listed as an author. The correctness of the proofs does not depend on trusting the AI: every theorem is checked by the Lean~4 kernel on
the empty axiom list. The fidelity of the definitions to the philosophical questions, and the literature check, do
depend on reading them, and the reader is invited to.

\begin{thebibliography}{99}\small
\bibitem{arnauld} A.~Arnauld. Fourth Objections. In R.~Descartes, \emph{Meditationes de prima philosophia}, 1641.
\bibitem{gupta} A.~Gupta, N.~Belnap. \emph{The Revision Theory of Truth}. MIT Press, 1993.
\bibitem{peirce} C.~S. Peirce. Questions concerning certain faculties claimed for man. \emph{Journal of
  Speculative Philosophy} 2 (1868) 103--114.
\bibitem{davidson} D.~Davidson. A coherence theory of truth and knowledge. In E.~LePore (ed.), \emph{Truth and
  Interpretation}, Blackwell, 1986.
\bibitem{hilbert} D.~Hilbert. Über das Unendliche. \emph{Mathematische Annalen} 95 (1926) 161--190.
\bibitem{fages} F.~Fages. Consistency of Clark's completion and existence of stable models. \emph{Journal of Methods of
  Logic in Computer Science} 1 (1994) 51--60.
\bibitem{hegel} G.~W.~F. Hegel. \emph{Phänomenologie des Geistes}. 1807.
\bibitem{albert} H.~Albert. \emph{Traktat über kritische Vernunft}. Mohr, Tübingen, 1968.
\bibitem{field} H.~Field. \emph{Science Without Numbers}. Princeton University Press, 1980.
\bibitem{gaifman} H.~Gaifman. Pointers to truth. \emph{Journal of Philosophy} 89 (1992) 223--261.
\bibitem{herzberger} H.~Herzberger. Paradoxes of grounding in semantics. \emph{Journal of Philosophy} 67 (1970) 145--167.
\bibitem{leitgeb} H.~Leitgeb. What truth depends on. \emph{Journal of Philosophical Logic} 34 (2005) 155--192.
\bibitem{kant} I.~Kant. \emph{Kritik der reinen Vernunft}. 1781/1787.
\bibitem{barnes} J.~Barnes. \emph{The Toils of Scepticism}. Cambridge University Press, 1990.
\bibitem{barnes82} J.~Barnes. The beliefs of a Pyrrhonist. \emph{Proceedings of the Cambridge Philological Society} 28 (1982) 1--29.
\bibitem{hamkins} J.~D. Hamkins. The set-theoretic multiverse. \emph{Review of Symbolic Logic} 5 (2012) 416--449.
\bibitem{fichte} J.~G. Fichte. \emph{Grundlage der gesammten Wissenschaftslehre}. 1794.
\bibitem{hintikka} J.~Hintikka. Cogito, ergo sum: inference or performance? \emph{Philosophical Review} 71 (1962) 3--32.
\bibitem{godel} K.~Gödel. The consistency of the axiom of choice and of the generalized continuum hypothesis.
  \emph{PNAS} 24 (1938) 556--557.
\bibitem{popper} K.~Popper. \emph{Logik der Forschung}. Springer, 1935.
\bibitem{bonjour} L.~BonJour. \emph{The Structure of Empirical Knowledge}. Harvard University Press, 1985.
\bibitem{wittgenstein} L.~Wittgenstein. \emph{On Certainty}. Blackwell, Oxford, 1969.
\bibitem{burnyeat} M.~Burnyeat. Can the sceptic live his scepticism? In M.~Schofield et al. (eds.), \emph{Doubt and
  Dogmatism}, Oxford, 1980.
\bibitem{frede} M.~Frede. Des Skeptikers Meinungen. \emph{Neue Hefte für Philosophie} 15/16 (1979) 102--129.
\bibitem{rathjen} M.~Rathjen. Indefiniteness in semi-intuitionistic set theories: on a conjecture of Feferman. \emph{Journal of Symbolic Logic} 81 (2016) 742--754.
\bibitem{richardson} M.~Richardson. Solutions of irreflexive relations. \emph{Annals of Mathematics} 58 (1953) 573--590.
\bibitem{walicki22} M.~Walicki, S.~Dyrkolbotn. Paraconsistent resolution. \emph{Australasian Journal of Logic} 19(3)
  (2022) 96--123.
\bibitem{walicki09} M.~Walicki. Reference, paradoxes and truth. \emph{Synthese} 171 (2009) 195--226.
\bibitem{cohen} P.~J. Cohen. The independence of the continuum hypothesis, I and II. \emph{PNAS} 50 (1963)
  1143--1148; 51 (1964) 105--110.
\bibitem{klein} P.~Klein. Human knowledge and the infinite regress of reasons. \emph{Philosophical Perspectives} 13
  (1999) 297--325.
\bibitem{dung} P.~M. Dung. On the acceptability of arguments. \emph{Artificial Intelligence} 77 (1995) 321--357.
\bibitem{fogelin} R.~Fogelin. \emph{Pyrrhonian Reflections on Knowledge and Justification}. Oxford University Press,
  1994.
\bibitem{cook} R.~T. Cook. Patterns of paradox. \emph{Journal of Symbolic Logic} 69 (2004) 767--774.
\bibitem{berker} S.~Berker. Coherentism via graphs. \emph{Philosophical Issues} 25 (2015) 322--352.
\bibitem{feferman} S.~Feferman. Is the continuum hypothesis a definite mathematical problem? EFI project draft, 2011.
\bibitem{kripke} S.~Kripke. Outline of a theory of truth. \emph{Journal of Philosophy} 72 (1975) 690--716.
\bibitem{shapiro} S.~Shapiro. Conservativeness and incompleteness. \emph{Journal of Philosophy} 80 (1983) 521--531.
\bibitem{yablo82} S.~Yablo. Grounding, dependence, and paradox. \emph{Journal of Philosophical Logic} 11 (1982) 117--137.
\bibitem{yablo} S.~Yablo. Paradox without self-reference. \emph{Analysis} 53 (1993) 251--252.
\bibitem{sextus} Sextus Empiricus. \emph{Outlines of Pyrrhonism}, transl. R.~G. Bury. Loeb Classical Library, 1933.
\bibitem{miao} T.~Miao. Reducing causality to functions with structural models. arXiv:2307.07524, 2023.
\bibitem{day} V.~Day, C.~Athos. The amphiboly in the third horn. Zenodo, 2026, doi:10.5281/zenodo.18888781.
\bibitem{reznikForms} V.~Reznik. VR-Forms: An operational two-register apparatus over VR-Sets, version 1.0.3.
  Zenodo, 2026, doi:10.5281/zenodo.23647163.
\bibitem{reznikZTL} V.~Reznik. Zero-Trust Logic, version 2.1.2. Zenodo, 2026, doi:10.5281/zenodo.23019513.
\bibitem{sellars} W.~Sellars. Empiricism and the philosophy of mind. \emph{Minnesota Studies in the Philosophy of
  Science} 1 (1956) 253--329.
\bibitem{quine} W.~V.~O. Quine. Two dogmas of empiricism. \emph{Philosophical Review} 60 (1951) 20--43.
\end{thebibliography}

\end{document}
